% ============================================================
% 全局宏包与样式定义（由 main.tex 在前言区 \input）
% 所有章节共享的包、颜色、环境和命令统一定义于此
% ============================================================

% --- 数学 ---
\usepackage{amsmath}
\usepackage{amssymb}

% --- 脚注每页重新编号 ---
\usepackage[perpage]{footmisc}

% --- 超链接（ctexbook 兼容模式）---
\usepackage[hidelinks,bookmarks=true]{hyperref}

% --- 引用 ---
\usepackage{csquotes}

% --- TikZ ---
\usepackage{tikz}
\usetikzlibrary{arrows,shapes,positioning,calc}

% --- 代码块 ---
\usepackage{listings}
\definecolor{codebg}{RGB}{245,245,250}
\definecolor{codeframe}{RGB}{180,180,200}
\lstdefinestyle{codeblock}{
    backgroundcolor=\color{codebg},
    basicstyle=\small\ttfamily,
    breaklines=true,
    frame=single,
    framerule=0.5pt,
    rulecolor=\color{codeframe},
    framesep=8pt,
    aboveskip=8pt,
    belowskip=8pt,
    columns=fullflexible,
    keepspaces=true,
    showstringspaces=false,
}
\lstset{style=codeblock}

% ============================================================
% 正式表格样式（三线表）
% ============================================================
\newcolumntype{Y}{>{\RaggedRight\arraybackslash}X}

\newenvironment{formalTable}[1][htbp]{%
    \begin{table}[#1]
    \centering
    \small
    \renewcommand{\arraystretch}{1.35}
}{%
    \end{table}
}

\newenvironment{formalLongTable}{%
    \begingroup
    \small
    \renewcommand{\arraystretch}{1.35}
    \setlength{\LTleft}{\fill}
    \setlength{\LTright}{\fill}
}{%
    \endgroup
}

% 非浮动表格（用于 tcolorbox 等不可含 float 的环境内部）
\newenvironment{formalInlineTable}{%
    \begin{center}
    \small
    \renewcommand{\arraystretch}{1.35}
}{%
    \end{center}
}

% 长表格统一跨页表头（在 longtable 内使用）
% #1 = 列数, #2 = 表头行内容
\newcommand{\formallongtablehead}[2]{%
    #2 \\
    \midrule
    \endfirsthead
    \multicolumn{#1}{l}{\small\sffamily （续表）}\\
    \addlinespace
    #2 \\
    \midrule
    \endhead
    \midrule
    \multicolumn{#1}{r}{\small 接下页}\\
    \endfoot
    \bottomrule
    \endlastfoot
}

% ============================================================
% tcolorbox 样式定义（黑白打印适配版）
% ============================================================
\tcbset{
    % 引用框样式
    textbookquote/.style={
        colback=gray!3,        % 极浅灰背景
        colframe=gray!40,      % 浅灰边框
        boxrule=0.5pt, arc=3pt,
        left=6pt, right=6pt,
        top=4pt, bottom=4pt,
    },    
    % 知识要点框样式（全书统一）
    skillbox/.style={
        enhanced,
        breakable,
        colback=gray!4,        % 浅灰背景
        colframe=gray!70,      % 深灰边框
        boxrule=0.7pt, arc=2mm,
        left=8pt, right=8pt,
        top=7pt, bottom=7pt,
        fonttitle=\bfseries,
    },
    % 每节开头的知识要点框样式（全书统一）
    sectionkeypointsbox/.style={
        skillbox,
        title={本节知识要点},
        colback=gray!3,
        colframe=gray!70,
        colbacktitle=gray!80,
        coltitle=white,
        before skip=4pt,
        after skip=4pt,
    },
    % AI 回答输出框样式
    aiboxsty/.style={
        skillbox, 
        colback=gray!4,        % 浅灰背景
        colframe=gray!60,      % 中深灰边框
        fontupper=\small, breakable,
        title={{\bfseries AI 输出}},
        coltitle=black,        % 标题使用纯黑以保证清晰
        before skip=4pt, after skip=4pt,
    },
    % 用户提示词输入框样式
    userboxsty/.style={
        skillbox, 
        colback=gray!4,        % 浅灰背景
        colframe=gray!50,      % 中灰边框
        fontupper=\small, breakable,
        title={{\bfseries AI 提问输入}},
        coltitle=black,
        before skip=4pt, after skip=4pt,
    },
    % % 科研提示词输入框样式（这部分与{userbox}重复了，不用这个）
    % researcherboxsty/.style={
    %     skillbox, 
    %     colback=gray!4,        % 浅灰背景
    %     colframe=gray!50,      % 中灰边框
    %     fontupper=\small, breakable,
    %     title={{\bfseries 输入}},
    %     coltitle=black,
    %     before skip=4pt, after skip=4pt,
    % },
    % 审视重点分析框样式
    critiqueboxsty/.style={
        skillbox, 
        colback=gray!3,        % 极浅灰背景
        colframe=gray!60,      % 中深灰边框
        fontupper=\small, breakable,
        title={{\bfseries 审核重点}},
        coltitle=black,
        before skip=4pt, after skip=4pt,
    },
    % 常见误区框样式
     mistakeboxsty/.style={
        skillbox, 
        colback=gray!3,        % 极浅灰背景
        colframe=gray!60,      % 中深灰边框
        fontupper=\small, breakable,
        title={{\bfseries 常见误区}},
        coltitle=black,
        before skip=4pt, after skip=4pt,
    },
    % 提示词模板展示框样式
    promptbox/.style={
        enhanced, breakable,
        colback=gray!4,        % 浅灰背景
        colframe=gray!60,      % 中深灰边框
        fonttitle=\bfseries,
        title=提示词,
        arc=1.5mm, boxrule=0.6pt,
        left=8pt, right=8pt, top=6pt, bottom=6pt,
        before skip=8pt, after skip=8pt,
    },
    % 操作步骤框样式
    stepbox/.style={
        enhanced, breakable,
        colback=gray!3,        % 极浅灰背景
        colframe=gray!55,      % 中灰边框
        fonttitle=\bfseries,
        arc=1.5mm, boxrule=0.6pt,
        left=8pt, right=8pt, top=6pt, bottom=6pt,
        before skip=8pt, after skip=8pt,
    },
    % 文件内容展示框样式
    filebox/.style={
        enhanced, breakable,
        colback=gray!4,        % 浅灰背景
        colframe=gray!65,      % 中深灰边框
        fonttitle=\bfseries,
        arc=1.5mm, boxrule=0.6pt,
        left=8pt, right=8pt, top=6pt, bottom=6pt,
        before skip=8pt, after skip=8pt,
    },
    % 趣味知识框样式
    funboxsty/.style={
        enhanced, breakable,
        colback=gray!4,        % 浅灰背景
        colframe=gray!60,      % 中深灰边框
        fonttitle=\bfseries,
        title=知识拓展,
        arc=2mm, boxrule=0.5pt,
        left=8pt, right=8pt, top=6pt, bottom=6pt,
        before skip=8pt, after skip=8pt,
    },
    % 第4章专用框体样式（通过边框粗细和灰度区分）
    ch4box-keypoints/.style={skillbox, colback=gray!3, colframe=gray!70, colbacktitle=gray!80, coltitle=white, fonttitle=\bfseries},
    ch4box-knowledge/.style={skillbox, colback=gray!4, colframe=gray!65, colbacktitle=gray!70, coltitle=white, fonttitle=\bfseries},
    ch4box-teacher/.style={skillbox, colback=gray!3, colframe=gray!75, colbacktitle=gray!85, coltitle=white, fonttitle=\bfseries},
    ch4box-prompt/.style={skillbox, colback=gray!4, colframe=gray!60, colbacktitle=gray!65, coltitle=white, fonttitle=\bfseries},
    ch4box-outputA/.style={skillbox, colback=gray!3, colframe=gray!70, colbacktitle=gray!75, coltitle=white, fonttitle=\bfseries},
    ch4box-outputB/.style={skillbox, colback=gray!4, colframe=gray!65, colbacktitle=gray!70, coltitle=white, fonttitle=\bfseries},
    ch4box-reminder/.style={skillbox, colback=gray!3, colframe=gray!60, colbacktitle=gray!75, coltitle=white, fonttitle=\bfseries},
}

% ============================================================
% 环境定义（保持不变）
% ============================================================

\newenvironment{lightquote}{%
    \begin{tcolorbox}[textbookquote]\small}{%
    \end{tcolorbox}}

\newenvironment{aibox}{%
    \begin{tcolorbox}[aiboxsty]}{%
    \end{tcolorbox}}

\newenvironment{userbox}{%
    \begin{tcolorbox}[userboxsty]}{%
    \end{tcolorbox}}

% \newenvironment{researcherbox}{%
%     \begin{tcolorbox}[researcherboxsty]}{%
%     \end{tcolorbox}}

\newenvironment{critiquebox}{%
    \begin{tcolorbox}[critiqueboxsty]}{%
    \end{tcolorbox}}

\newenvironment{mistakebox}{%
    \begin{tcolorbox}[mistakeboxsty]}{%
    \end{tcolorbox}}



\newenvironment{funbox}{%
    \begin{tcolorbox}[funboxsty]\small}{%
    \end{tcolorbox}}

\newenvironment{warnbox}{%
    \begin{tcolorbox}[warnbox]\small}{%
    \end{tcolorbox}}

\newenvironment{promptbox}{%
    \begin{tcolorbox}[promptbox]}{%
    \end{tcolorbox}}

\newenvironment{knowledgebox}{%
    \begin{tcolorbox}[skillbox, title=知识拓展]}{%
    \end{tcolorbox}}

\newenvironment{sectionkeypoints}{%
    \begin{tcolorbox}[sectionkeypointsbox]\small
    \begin{itemize}[leftmargin=*, itemsep=2pt]}{%
    \end{itemize}
    \end{tcolorbox}}
% % ============================================================
% % tcolorbox 样式定义（有颜色，仅定义样式，不创建环境）
% % ============================================================
% \tcbset{
%     % 引用框样式
%     textbookquote/.style={
%         colback=blue!3!white,
%         colframe=blue!30!white,
%         boxrule=0.5pt, arc=3pt,
%         left=6pt, right=6pt,
%         top=4pt, bottom=4pt,
%     },    
%     % 知识要点框样式（全书统一）
%     skillbox/.style={
%         enhanced,
%         breakable,
%         colback=blue!4, colframe=blue!55!black,
%         boxrule=0.7pt, arc=2mm,
%         left=8pt, right=8pt,
%         top=7pt, bottom=7pt,
%         fonttitle=\bfseries,
%     },
%     % 每节开头的知识要点框样式（全书统一）
%     sectionkeypointsbox/.style={
%         skillbox,
%         title={本节知识要点},
%         colback=blue!3,
%         colframe=blue!55!black,
%         colbacktitle=blue!75!black,
%         coltitle=white,
%         before skip=4pt,
%         after skip=4pt,
%     },
%     % AI 回答输出框样式
%     aiboxsty/.style={
%         skillbox, colback=teal!4, colframe=teal!50!black,
%         fontupper=\small, breakable,
%         title={{\bfseries AI 输出}},
%         coltitle=teal!50!black!70,
%         before skip=4pt, after skip=4pt,
%     },
%     % 用户提示词输入框样式
%     userboxsty/.style={
%         skillbox, colback=blue!4, colframe=blue!50!black!30,
%         fontupper=\small, breakable,
%         title={{\bfseries AI 提问输入}},
%         coltitle=blue!50!black!70,
%         before skip=4pt, after skip=4pt,
%     },
%     % 科研提示词输入框样式
%     researcherboxsty/.style={
%         skillbox, colback=blue!4, colframe=blue!50!black!30,
%         fontupper=\small, breakable,
%         title={{\bfseries 输入}},
%         coltitle=blue!50!black!70,
%         before skip=4pt, after skip=4pt,
%     },
%     % 审视/批判分析框样式
%     critiqueboxsty/.style={
%         skillbox, colback=cyan!3, colframe=cyan!45!black,
%         fontupper=\footnotesize, breakable,
%         title={{\bfseries 审视}},
%         coltitle=cyan!50!black!70,
%         before skip=4pt, after skip=4pt,
%     },
%     % 提示词模板展示框样式
%     promptbox/.style={
%         enhanced, breakable,
%         colback=gray!4, colframe=gray!55,
%         fonttitle=\bfseries,
%         title=提示词或示例内容,
%         arc=1.5mm, boxrule=0.6pt,
%         left=8pt, right=8pt, top=6pt, bottom=6pt,
%         before skip=8pt, after skip=8pt,
%     },
%     % 操作步骤框样式
%     stepbox/.style={
%         enhanced, breakable,
%         colback=blue!3, colframe=blue!45!black,
%         fonttitle=\bfseries,
%         arc=1.5mm, boxrule=0.6pt,
%         left=8pt, right=8pt, top=6pt, bottom=6pt,
%         before skip=8pt, after skip=8pt,
%     },
%     % 文件内容展示框样式
%     filebox/.style={
%         enhanced, breakable,
%         colback=gray!4, colframe=gray!60,
%         fonttitle=\bfseries,
%         arc=1.5mm, boxrule=0.6pt,
%         left=8pt, right=8pt, top=6pt, bottom=6pt,
%         before skip=8pt, after skip=8pt,
%     },
%     % 警告/注意框样式
%     warnbox/.style={
%         enhanced, breakable,
%         colback=gray!5, colframe=gray!65,
%         fonttitle=\bfseries,
%         title=注意,
%         arc=1.5mm, boxrule=0.6pt,
%         left=8pt, right=8pt, top=6pt, bottom=6pt,
%         before skip=8pt, after skip=8pt,
%     },
%     % 趣味知识框样式
%     funboxsty/.style={
%         enhanced, breakable,
%         colback=yellow!6, colframe=orange!40!black,
%         fonttitle=\bfseries,
%         title=知识拓展,
%         arc=2mm, boxrule=0.5pt,
%         left=8pt, right=8pt, top=6pt, bottom=6pt,
%         before skip=8pt, after skip=8pt,
%     },
%     % 第4章专用框体样式
%     ch4box-keypoints/.style={skillbox, colback=blue!3, colframe=blue!55!black, colbacktitle=blue!75!black, coltitle=white, fonttitle=\bfseries},
%     ch4box-knowledge/.style={skillbox, colback=teal!4, colframe=teal!55!black, colbacktitle=teal!70!black, coltitle=white, fonttitle=\bfseries},
%     ch4box-teacher/.style={skillbox, colback=orange!6, colframe=orange!70!black, colbacktitle=orange!75!black, coltitle=white, fonttitle=\bfseries},
%     ch4box-prompt/.style={skillbox, colback=purple!4, colframe=purple!60!black, colbacktitle=purple!70!black, coltitle=white, fonttitle=\bfseries},
%     ch4box-outputA/.style={skillbox, colback=red!3, colframe=red!55!black, colbacktitle=red!65!black, coltitle=white, fonttitle=\bfseries},
%     ch4box-outputB/.style={skillbox, colback=green!5, colframe=green!50!black, colbacktitle=green!60!black, coltitle=white, fonttitle=\bfseries},
%     ch4box-reminder/.style={skillbox, colback=cyan!4, colframe=cyan!55!black, colbacktitle=cyan!70!black, coltitle=white, fonttitle=\bfseries},
% }

% % ============================================================
% % 环境定义（使用 newenvironment 包装，比 \newtcolorbox 更稳健）
% % ============================================================

% \newenvironment{lightquote}{%
%     \begin{tcolorbox}[textbookquote]\small}{%
%     \end{tcolorbox}}

% \newenvironment{aibox}{%
%     \begin{tcolorbox}[aiboxsty]}{%
%     \end{tcolorbox}}

% \newenvironment{userbox}{%
%     \begin{tcolorbox}[userboxsty]}{%
%     \end{tcolorbox}}

% \newenvironment{researcherbox}{%
%     \begin{tcolorbox}[researcherboxsty]}{%
%     \end{tcolorbox}}

% \newenvironment{critiquebox}{%
%     \begin{tcolorbox}[critiqueboxsty]}{%
%     \end{tcolorbox}}

% \newenvironment{funbox}{%
%     \begin{tcolorbox}[funboxsty]\small}{%
%     \end{tcolorbox}}

% \newenvironment{knowledgebox}{%
%     \begin{tcolorbox}[skillbox, title=知识拓展]}{%
%     \end{tcolorbox}}

% ============================================================
% 通用命令定义
% ============================================================

\providecommand{\minorheading}[1]{%
    \par\medskip
    \noindent\textbf{#1}\par\vspace{0.25em}
}

\providecommand{\manualfigcaption}[1]{%
    \par\vspace{0.35em}{\centering\small #1\par}\vspace{0.25em}%
}

\providecommand{\BookFigureImage}[1]{%
    \includegraphics[width=0.93\textwidth,height=0.42\textheight,keepaspectratio]{#1}%
}

\providecommand{\caselabel}[2]{{\small\sffamily\color{#1}#2}}
\providecommand{\casestar}{$\bigstar$}
